%%% FIELD proof-fuzzy-membership
Put \(B_0<\infty\) for the fixed overlap multiplicity of the profiles.  Scaling the fixed \(C^\infty\) bumps in the profile construction gives, for every fixed \(t>0\),
\[
 \|F_\lambda\|_{\mathcal H^t}\le C_t\ell^{-t},
 \qquad
 \|K_h\|_{\mathcal H^t}\le C_t h^{-t},
 \qquad
 \|F_\lambda\|_\infty\le C_0,
 \qquad 0\le K_h\le1 .
\]
Indeed, differentiation through order \(\lfloor t\rfloor\) contributes the corresponding inverse scale and the last derivative has the required Hölder increment; at each point at most \(B_0\) profiles occur.  These bounds are uniform in \(k\) and \(\lambda\).  Since \(a\le1/8\),
\[
 D_{\mathrm{fuz}}\ge \frac14-\frac1{64}=\frac{15}{64}.
\]
The quotient rule through order \(\lfloor t\rfloor\), followed by the same Hölder-increment bound, therefore yields
\[
 G_h:=\frac{K_h}{D_{\mathrm{fuz}}},
 \qquad
 \|G_h\|_{\mathcal H^t}\le C_t h^{-t}.
\]
All divisions used below are consequently legitimate.

The propensity satisfies
\[
 \|\pi_\lambda-1/2\|_{\mathcal H^\alpha}
 \le C a\ell^{-\alpha}\le Cc.
\]
Moreover \(q=\alpha+\beta<\gamma\) implies \(\gamma>\beta\), and
\[
 \tau_s=-sabG_h,
 \qquad
 \|\tau_s\|_{\mathcal H^\gamma}
 \le C ab h^{-\gamma}\le Cc,
 \qquad
 \|\tau_s\|_{\mathcal H^\beta}
 \le C ab h^{-\beta}\le Cc h^{\gamma-\beta}\le Cc.
\]
Using \(\mu_{0,s,\lambda}=1/2+sbF_\lambda-(1/2+aF_\lambda)\tau_s\), the scaled product rule and \(\ell\le h\) give
\[
 \begin{aligned}
 \|bF_\lambda\|_{\mathcal H^\beta}
   &\le Cb\ell^{-\beta}\le Cc,\\
 \|aF_\lambda\tau_s\|_{\mathcal H^\beta}
   &=\|a^2bF_\lambda G_h\|_{\mathcal H^\beta}
     \le Ca^2b\ell^{-\beta}
     \le Cb\ell^{-\beta}\le Cc.
 \end{aligned}
\]
Here the product estimate follows directly by distributing each derivative or Hölder increment between a scale-\(\ell\) factor and a scale-\(h\) factor; because \(\ell\le h\), the largest inverse scale is \(\ell^{-\beta}\).  Thus \(\mu_{0,s,\lambda}\in\mathcal H^\beta(L)\), \(\pi_\lambda\in\mathcal H^\alpha(L)\), and \(\tau_s\in\mathcal H^\gamma(L)\) once \(c>0\) is chosen sufficiently small in terms of the displayed fixed parameters and \(L\ge2\).

The same choice of \(c\) makes \(C_0a\le1/4\), \(C_0b\le1/4\), and
\[
 |\tau_s|\le \frac{ab}{15/64}.
\]
Consequently \(e_0\le\pi_\lambda\le1-e_0\), because \(e_0\le1/4\), and both \(\mu_{0,s,\lambda}\) and \(\mu_{1,s,\lambda}=\mu_{0,s,\lambda}+\tau_s\) lie in \((0,1)\).  The construction gives consistency, conditional exchangeability, and the exact uniform design.  Hence every \(P_{s,\lambda}\) belongs to \(\mathcal M_H\).  Since \(K_h(x_0)=1\),
\[
 \theta(P_{s,\lambda})=\tau_s(x_0)
 =-\frac{sab}{1/4-a^2}.
\]

For \(P_{\mathrm{const},\xi}\), the propensity and both arm means are constant, the effect is zero, and the same smallness choice puts the propensity in the overlap interval and the means in \((0,1)\).  All Hölder requirements therefore hold and \(\theta(P_{\mathrm{const},\xi})=0\).  Finally average the two observed constant-component laws.  Their treatment probability is \(1/2\), while
\[
 \begin{aligned}
 P(A=1,Y=1)&=\frac12\sum_{\xi\in\{-1,1\}}
 (1/2+a\xi)(1/2+b\xi)=\frac14+ab,\\
 P(A=0,Y=1)&=\frac12\sum_{\xi\in\{-1,1\}}
 (1/2-a\xi)(1/2+b\xi)=\frac14-ab.
 \end{aligned}
\]
Thus \(Q_{\mathrm{const}}\) has constant propensity \(1/2\), arm means \(1/2+2ab\) and \(1/2-2ab\), and effect \(4ab\).  Its stipulated causal realization satisfies consistency and conditional exchangeability, and its constant functions satisfy every Hölder and overlap condition.  Therefore \(Q_{\mathrm{const}}\in\mathcal M_H\) and \(\theta(Q_{\mathrm{const}})=4ab\), as claimed.

%%% FIELD proof-one-record-obstruction
The first equality is part of the construction of \(Q_{\mathrm{const}}\) as an observed-data law.  The preceding lemma proves that all three laws are in \(\mathcal M_H\), that the two components have target zero, and that \(Q_{\mathrm{const}}\) has target \(4ab\).

For the signed family, fix \(x\) and write \(F=F_\lambda(x)\), \(K=K_h(x)\), \(p=\pi_\lambda(x)\), and \(m=m_{s,\lambda}(x)\).  With \(U=2A-1\) and \(V=2Y-1\), the construction implies
\[
 \mathbb E(U\mid X=x)=2aF,
 \qquad
 \mathbb E(V\mid X=x)=2sbF.
\]
Because \(m=(1-p)\mu_0+p\mu_1\) and \(\mu_1-\mu_0=\tau_s\), direct expansion gives
\[
 \begin{aligned}
 \mathbb E(UV\mid X=x)
 &=4\{pm+p(1-p)\tau_s\}-2p-2m+1\\
 &=(2p-1)(2m-1)+4p(1-p)\tau_s\\
 &=4sabF^2+4(1/4-a^2F^2)\tau_s.
 \end{aligned}
\]
Let \(\mathbb E_\lambda\) denote uniform averaging over \(\{-1,1\}^k\).  Independence and symmetry of the coordinates give
\[
 \mathbb E_\lambda F_\lambda(x)=0,
 \qquad
 \mathbb E_\lambda F_\lambda(x)^2=K_h(x).
\]
Since \(\tau_s=-sabK_h/(1/4-a^2K_h)\) does not depend on \(\lambda\),
\[
 \mathbb E_\lambda\mathbb E(UV\mid X=x)
 =4sabK_h+4(1/4-a^2K_h)\tau_s=0.
\]
The law of a pair of signs is determined by its three moments through
\[
 \Pr(U=u,V=w\mid X=x)
 =\frac14\{1+u\mathbb E U+w\mathbb E V+uw\mathbb E(UV)\}.
\]
After averaging \(\lambda\), all three nonconstant moments vanish.  Hence the conditional observed mark law is the uniform law on \(\{-1,1\}^2\), independently of \(s\).  The \(X\)-marginal is the same uniform law for both signs, so
\[
 2^{-k}\sum_\lambda P_{+,\lambda}
 =2^{-k}\sum_\lambda P_{-,\lambda}
\]
on one observed record.

The constant-law identity places \(Q_{\mathrm{const}}\) simultaneously in the one-record convex hull of target-zero laws and in the target-\(4ab\) fiber.  The displayed signed identity gives the analogous intersection of the two signed prior convex hulls.  Finally, for any bounded measurable one-record score \(g\), equality of the signed mixtures gives
\[
 2^{-k}\sum_\lambda\mathbb E_{P_{+,\lambda}}g(O)
 -2^{-k}\sum_\lambda\mathbb E_{P_{-,\lambda}}g(O)=0.
\]
Summing such terms over records preserves zero margin, proving the assertion for every bounded additive one-record score.

%%% FIELD proof-pair-mean-sensitivity
On \(C_j\), the profile properties give \(\zeta_h=1\) and \(\psi_j=1\).  Since the nonnegative squares sum to one there, every other \(\psi_l\) vanishes.  Thus \(F_\lambda=\lambda_j\) on \(C_j\), and
\[
 \mathbb E(U\mid X\in C_j)=2a\lambda_j,
 \qquad
 \mathbb E(V\mid X\in C_j)=2sb\lambda_j.
\]
For two independent records, conditional independence across record indices therefore gives
\[
 \begin{aligned}
 \mathbb E(Z_{12}\mid X_1,X_2\in C_j)
 &=\frac12\{\mathbb E U_1\mathbb E V_2+\mathbb E V_1\mathbb E U_2\}\\
 &=4sab\lambda_j^2=4sab.
 \end{aligned}
\]
Conditional on \(N_j=m\ge2\), the records in \(C_j\) are exchangeable draws from the law conditional on \(X\in C_j\).  Every ordered summand consequently has mean \(4sab\), so
\[
 \mathbb E(B_j\mid N_j)=4sab\mathbf1\{N_j\ge2\}.
\]
Uniform design gives \(N_j\sim\operatorname{Bin}(n,v)\); hence
\[
 \mathbb E S_{\mathrm{pair}}
 =\sum_{j=1}^k4sab\Pr(N_j\ge2)=4sabkp_2.
\]

It remains to prove the all-dataset sensitivity assertion.  Every \(Z_{ir}\) lies in \([-1,1]\).  If a cell contains \(m\ge2\) records, deleting one record changes its ordered-pair average by at most \(4/m\le2\): the retained ordered pairs have total weight change \(2/m\), and the deleted incident ordered pairs have total weight \(2/m\), with summands bounded by one.  The boundary change from two records to one has size at most one.  The same bound follows for insertion by reversing the operation.  Replacing one complete record can affect only its old and new core cells.  If these are different, each cell changes by at most two; if they coincide, deletion followed by insertion gives the same total bound.  Therefore
\[
 |S_{\mathrm{pair}}(D_n)-S_{\mathrm{pair}}(D_n')|\le4
 \qquad\text{whenever }D_n\sim_1D_n',
\]
and taking the supremum proves \(\Delta_{\mathrm{pair}}\le4\).

%%% FIELD proof-pair-composite-test
We first record the occupancy estimate needed for concentration.  Conditional on all covariates, let
\[
 I_j=\mathbf1\{N_j\ge2\},
 \qquad M=\sum_{j=1}^k I_j.
\]
The multinomial counts \(N_1,\ldots,N_k\) are negatively associated.  Here this fact follows directly from the sampling construction: the category-indicator vector of one record is one-hot and hence negatively associated; conditioning and the covariance decomposition show inductively that an independent union of negatively associated families remains negatively associated; coordinatewise summation over records on disjoint coordinate blocks preserves the defining covariance inequality.  Applying these two closure steps to the \(n\) independent one-hot vectors gives the assertion.  Replacing increasing functions by their negatives gives the same product inequality for nonnegative decreasing one-coordinate functions.  Consequently, for every \(t\ge0\),
\[
 \begin{aligned}
 \mathbb E e^{-tM}
 &\le\prod_{j=1}^k\mathbb E e^{-tI_j}
   =\{1-p_2(1-e^{-t})\}^k\\
 &\le\exp\{-kp_2(1-e^{-t})\}.
 \end{aligned}
\]
For \(0\le t\le1\), \(1-e^{-t}\ge t/2\), and thus
\[
 \mathbb E e^{-tM}\le e^{-kp_2t/2}. \tag{1}
\]
This derivation uses only independent sampling and the public disjoint cells; it introduces no nuisance-cardinality factor.

Conditional on the covariates, the active \(B_j\)'s use disjoint sets of independent marks and are therefore independent.  By the preceding lemma,
\[
 \mathbb E(sB_j\mid X_1,\ldots,X_n)=4abI_j,
 \qquad -1\le sB_j\le1.
\]
The elementary bounded-variable exponential estimate
\(\mathbb E\exp\{t(W-\mathbb EW)\}\le\exp(t^2/2)\) for \(W\in[-1,1]\), obtained by convexity of the exponential followed by optimization in \(t\), and Chernoff's inequality yield
\[
 \Pr\{sS_{\mathrm{pair}}\le2abM\mid X_1,\ldots,X_n\}
 \le \exp(-cMa^2b^2). \tag{2}
\]
If \(s(S_{\mathrm{pair}}+W)\le0\) while the complementary event in \(2\) holds, then \(sW\le-2abM\).  Since \(sW\) again has the \(\operatorname{Laplace}(4/\epsilon)\) law,
\[
 \Pr\{sW\le-2abM\mid X_1,\ldots,X_n\}
 =\frac12\exp(-\epsilon abM/2). \tag{3}
\]
Average \(2\)--\(3\) over the covariates and apply \(1\), first with \(t=ca^2b^2\) and then with \(t=\epsilon ab/2\).  Both are at most one because \(a,b\le1/8\) and \(\epsilon\le1\).  After changing numerical constants,
\[
 \sup_{s,\lambda}P_{s,\lambda}^{\otimes n}
 \{\phi_{\mathrm{pair}}\ne s\}
 \le C e^{-ckp_2a^2b^2}
      +C e^{-c\epsilon kp_2ab}. \tag{4}
\]

By the sensitivity lemma, the real-valued query \(S_{\mathrm{pair}}\) has replacement sensitivity at most four.  For neighboring datasets and every released real value \(z\), the two Laplace densities obey
\[
 \frac{\exp\{-\epsilon|z-S_{\mathrm{pair}}(D_n)|/4\}}
      {\exp\{-\epsilon|z-S_{\mathrm{pair}}(D_n')|/4\}}
 \le \exp\{\epsilon|S_{\mathrm{pair}}(D_n)-S_{\mathrm{pair}}(D_n')|/4\}
 \le e^\epsilon.
\]
Integration proves pure \(\epsilon\)-differential privacy of the noisy query.  The sign, fair tie-break, and scalar rescaling are postprocessings, so both \(\phi_{\mathrm{pair}}\) and \(\widehat\theta_{\mathrm{pair}}\) are pure \(\epsilon\)-differentially private.

Finally, fuzzy-family membership gives
\[
 \theta(P_{s,\lambda})=-sA_0,
 \qquad A_0:=\frac{ab}{1/4-a^2},
 \qquad \widehat\theta_{\mathrm{pair}}=-\phi_{\mathrm{pair}}A_0.
\]
The loss is zero when \(\phi_{\mathrm{pair}}=s\) and equals \(2A_0\) otherwise.  Multiplying \(4\) by \(2A_0\) proves the displayed risk bound.  No step selected among the \(2^k\) nuisance signs, so neither exponent contains a nuisance-sign logarithm.

%%% FIELD proof-two-record-exact-audit
For \(z_l=(j_l,u_l,w_l)\), expand the two factors in the definition and average the independent Rademacher coordinates of \(\lambda\).  Terms with a single \(\lambda_j\), or with two different coordinates, vanish.  Therefore
\[
 \overline\Pi_{s,k}^{\,2}(z_1,z_2)
 =\frac1{16k^2}\left[1+\mathbf1\{j_1=j_2\}
 \{4a^2u_1u_2+4b^2w_1w_2+4sab(u_1w_2+w_1u_2)\}\right]. \tag{1}
\]
The labels are independently uniform, so \(\Pr(\mathcal E_{\mathrm{coll}})=1/k\).  Subtracting \(1\) at the two signs gives
\[
 \overline\Pi_{+,k}^{\,2}(z_1,z_2)
 -\overline\Pi_{-,k}^{\,2}(z_1,z_2)
 =\frac{ab}{k^2}T_{2,k}(z_1,z_2). \tag{2}
\]
For a fixed common label, \(T_{2,k}\) has absolute value one on eight of the sixteen sign quadruples and is zero on the other eight.  Hence
\[
 \operatorname{TV}(\overline\Pi_{+,k}^{\,2},\overline\Pi_{-,k}^{\,2})
 =\frac12\frac{ab}{k^2}\,k\,8=\frac{4ab}{k}. \tag{3}
\]
Equation \(2\) also shows that the sign of \(T_{2,k}\), with a fair tie-break at zero, is a Bayes rule for equal prior weights.  Thus
\[
 \Pr(\sigma_{2,k}\ne s)=\frac{1-4ab/k}{2},
 \qquad
 \mathbb E(s\sigma_{2,k})=\frac{4ab}{k}. \tag{4}
\]
Write \(\rho=(e^\epsilon-1)/(e^\epsilon+1)=\tanh(\epsilon/2)\).  Randomized response satisfies \(\mathbb E(RR_{2,k}\mid\sigma_{2,k})=\rho\sigma_{2,k}\).  Averaging \(4\) therefore gives
\[
 \Pr(RR_{2,k}\ne s)
 =\frac12\{1-(4ab/k)\tanh(\epsilon/2)\}. \tag{5}
\]
The true and reported target values are respectively
\(-sab/(1/4-a^2)\) and \(-RR_{2,k}ab/(1/4-a^2)\).  The absolute loss is twice this amplitude exactly on the event in \(5\), which proves the asserted Bayes risk.

For an arbitrary input, including a tie input, each output probability of randomized response lies between \(1/(1+e^\epsilon)\) and \(e^\epsilon/(1+e^\epsilon)\).  Thus for every two inputs, and in particular every pair adjacent under \(\sim_{\mathrm q}\), each output likelihood ratio is at most \(e^\epsilon\).  Equality is attained on adjacent datasets: take
\[
 z_1=(j,1,1),\quad z_2=(j,1,1),\quad z_2'=(j,-1,-1).
\]
Then \(T_{2,k}(z_1,z_2)=1\) and \(T_{2,k}(z_1,z_2')=-1\), so the likelihood ratio for output \(+1\) is exactly \(e^\epsilon\).  The maximum neighboring log-likelihood ratio is consequently exactly \(\epsilon\).

%%% FIELD proof-pair-radius
We prove the proposition with the stated necessary boundary condition \(n\ge2\).  Since \(N_j\sim\operatorname{Bin}(n,v)\),
\[
 \binom n2v^2(1-v)^{n-2}
 \le p_2
 \le \mathbb E\binom{N_j}{2}=\binom n2v^2. \tag{1}
\]
The assumption \(nv\le c_0<1\) gives
\((1-v)^{n-2}\ge(1-v)^n\ge1-nv\ge1-c_0\), while \(n\ge2\) gives \(n(n-1)\asymp n^2\).  Thus
\[
 p_2\asymp n^2v^2,
 \qquad
 kp_2\asymp n^2\frac{h^d}{\ell^d}\ell^{2d}
 =n^2h^d\ell^d. \tag{2}
\]
All comparison constants depend only on the fixed geometric and occupancy constants.

The calibrations \(a\asymp\ell^\alpha\) and \(b\asymp\ell^\beta\), together with \(q=\alpha+\beta\), give
\[
 \delta\asymp ab\asymp\ell^q,
 \qquad
 \ell\asymp\delta^{1/q},
 \qquad
 h\asymp\delta^{1/\gamma}.
\]
Substitution in \(2\) yields
\[
 \begin{aligned}
 kp_2a^2b^2
 &\asymp n^2\delta^{,2+d/\gamma+d/q}
   =n^2\delta^{D_F+1},\\
 \epsilon kp_2ab
 &\asymp n^2\epsilon\delta^{,1+d/\gamma+d/q}
   =n^2\epsilon\delta^{D_F}. \tag{3}
 \end{aligned}
\]
Because \(r_{\mathrm{pair}}=n^{-2/(D_F+1)}\), the first quantity in \(3\) is at least a constant times \(C^{D_F+1}\) whenever \(\delta\ge Cr_{\mathrm{pair}}\).  Under \(n^{-1}\le\epsilon\le1\) and \(n\ge2\), \(n^2\epsilon>1\), so the cap in \(r_F\) is inactive and \(r_F=(n^2\epsilon)^{-1/D_F}\).  The second quantity in \(3\) is therefore at least a constant times \(C^{D_F}\) whenever \(\delta\ge Cr_F\).  Both exponents are bounded below for fixed \(C\) and diverge with \(C\), proving the fixed-error assertion at \(\max\{r_{\mathrm{pair}},r_F\}\).  The strict condition \(q<\gamma\) ensures
\(\ell/h\asymp\delta^{1/q-1/\gamma}\to0\) as \(\delta\downarrow0\), so this calibration respects the stipulated two-scale geometry.

For the final sanity check, set \(k=1\).  On \(C_1\), \(F_{(\lambda_1)}=\lambda_1\), whence
\[
 \mathbb E(U_i\mid X_i\in C_1)=2a\lambda_1,
 \qquad
 \mathbb E(V_i\mid X_i\in C_1)=2sb\lambda_1.
\]
Independence of the two records gives
\[
 \mathbb E\!\left[Z_{12}\mid X_1,X_2\in C_1\right]
 =\frac12\{(2a\lambda_1)(2sb\lambda_1)
 +(2sb\lambda_1)(2a\lambda_1)\}=4sab,
\]
as required.

%%% FIELD proof-shell-functional-risk
For neighboring datasets \(D_n\sim_1D_n'\), one-sensitivity gives
\[
 e^{-\epsilon/2}
 \le\frac{e^{-\epsilon Q_t(D_n)/2}}{e^{-\epsilon Q_t(D_n')/2}}
 \le e^{\epsilon/2}
\]
for every \(t\).  Summing these inequalities gives the reciprocal bound \(e^{\epsilon/2}\) for the ratio of the two normalizing constants.  Multiplying numerator and normalizer bounds shows that every output likelihood ratio of \(T_{\mathrm{shell}}\) is at most \(e^\epsilon\).  Thus the selector is pure \(\epsilon\)-differentially private under replacement adjacency.

Fix \(P\in\mathcal M_H\).  The maximum gap of \(\mathcal G_r\), including its final endpoint gap, is at most \(r\), so the nearest grid point satisfies
\[
 d_\Theta(t_P,\theta(P))\le r.
\]
For every dataset and every \(t\in\mathcal G_r\), retaining only the \(t_P\)-term in the selector denominator gives
\[
 \Pr(T_{\mathrm{shell}}=t\mid D_n)
 \le 1\wedge
 \exp\{-\epsilon(Q_t(D_n)-Q_{t_P}(D_n))/2\}. \tag{1}
\]
Every grid point with \(d_\Theta(t,t_P)<r\) contributes at most \(r\) in total to the conditional expectation of \(d_\Theta(T_{\mathrm{shell}},t_P)\).  On \(\mathcal S_m(P;r)\), \(d_\Theta(t,t_P)<(m+1)r\).  Therefore \(1\), Tonelli's theorem for the nonnegative finite sum, and the definition of \(\mathfrak W_r\) imply
\[
 \begin{aligned}
 \mathbb E_{P^{\otimes n}}d_\Theta(T_{\mathrm{shell}},t_P)
 &\le r+r\sum_{m\ge1}(m+1)
 \sum_{t\in\mathcal S_m(P;r)}
 \mathbb E_{P^{\otimes n}}
 \left[1\wedge e^{-\epsilon(Q_t-Q_{t_P})/2}\right]\\
 &\le (B+1)r.
 \end{aligned}
\]
The triangle inequality now gives
\[
 \mathbb E_{P^{\otimes n}}d_\Theta(T_{\mathrm{shell}},\theta(P))
 \le(B+2)r\le(B+3)r.
\]
Taking the supremum over \(P\) proves the claim.  The proof sums only the prescribed scalar distance shells, so neither the grid cardinality nor a nuisance-sieve cardinality enters the bound.

%%% FIELD partial-full-fiber
The finite signed family admits the explicit one-sensitive, log-free pair mechanism proved above: its error is bounded by
\[
 C\exp(-ckp_2a^2b^2)+C\exp(-c\epsilon kp_2ab),
\]
and the coupled privacy balance is \(r_F\).  Independently, any score family satisfying the stated one-sensitivity and weighted-shell certificate yields a pure \(\epsilon\)-differentially private estimator with risk at most \((B+3)r\).  What is not proved is that such scores exist uniformly over the entire Hölder class at \(r\) comparable to \(\max\{R_n^{\mathrm{NP}},r_G,r_F\}\).
